\documentclass{article}
\usepackage[margin=1in, a4paper]{geometry}
\usepackage{amsmath, amssymb, amsfonts}
\usepackage{graphicx}
\usepackage{xcolor}
\usepackage{listings}
\usepackage{booktabs}
\usepackage[utf8]{inputenc}
\usepackage[T1]{fontenc}
\usepackage{hyperref}
\hypersetup{colorlinks=true, urlcolor=blue, linkcolor=blue}
\usepackage{enumitem}
\usepackage{float}

\definecolor{codegreen}{rgb}{0,0.6,0}
\definecolor{codegray}{rgb}{0.5,0.5,0.5}
\definecolor{codepurple}{rgb}{0.58,0,0.82}
\definecolor{backcolour}{rgb}{0.99,0.99,0.99}
% Professional color palette for academic publishing
\definecolor{myblue}{cmyk}{0.8,0.4,0,0.1}
\definecolor{mygreen}{cmyk}{0.7,0,0.5,0.2}
\definecolor{myorange}{cmyk}{0,0.5,0.8,0.1}
\definecolor{mygray}{cmyk}{0,0,0,0.4}
\definecolor{arrowcolor}{cmyk}{0.9,0.5,0,0.3}
\definecolor{nodecolor}{cmyk}{0.6,0.2,0,0.05}
\definecolor{framecolor}{cmyk}{0.4,0.1,0,0.1}

\lstdefinestyle{mystyle}{
    backgroundcolor=\color{backcolour},   
    commentstyle=\color{codegreen},
    keywordstyle=\color{magenta},
    numberstyle=\tiny\color{codegray},
    stringstyle=\color{codepurple},
    basicstyle=\ttfamily\footnotesize,
    breakatwhitespace=false,         
    breaklines=true,                 
    captionpos=b,                    
    keepspaces=true,                 
    numbers=left,                    
    numbersep=5pt,                  
    showspaces=false,                
    showstringspaces=false,
    showtabs=false,                  
    tabsize=2
}
\lstset{style=mystyle}

\title{The Logistics \& Supply Chain Problem-Solving Playbook\\
Chapter 6: The Yard Location Naming Convention Problem}
\author{The Comprehensive Guide to Terminal Operations}
\date{}

\begin{document}

\maketitle

\section{The Yard Location Naming Convention Problem}

The efficient operation of container terminals, truck yards, and logistics facilities depends fundamentally on the ability to uniquely identify, locate, and reference every physical storage position within the facility. The yard location naming convention problem addresses the systematic design of naming schemes that enable rapid location identification, minimize human error, support scalable operations, and integrate seamlessly with warehouse management systems.

\subsection{The Scenario: Chaos in the Container Yard}

Pacific Gateway Terminal operates a 45-hectare container yard handling over 2.5 million TEUs annually. The facility contains 8,500 ground storage positions organized across 12 blocks, each with varying numbers of rows, bays, and tiers. Currently, the terminal uses an inconsistent naming system where different areas employ different conventions: Block A uses alphanumeric codes like ``A7-15-3'', Block B uses sequential numbers like ``2847'', while the newest expansion area uses GPS coordinates.

This naming chaos creates cascading operational problems. Yard truck drivers spend an average of 4.2 minutes per container searching for the correct storage position, resulting in 180 hours of lost productivity daily. The terminal information system contains over 800 duplicate location entries and 1,200 ``orphaned'' containers with invalid location codes. Container retrieval operations suffer from a 12\% error rate, with incorrect containers being pulled from yard positions, causing delays and additional handling costs.

The challenge intensifies during peak operations when temporary overflow areas are activated, requiring rapid integration of new storage positions into the existing naming system. Additionally, the terminal must accommodate different equipment types (reach stackers, rubber-tired gantry cranes, empty handlers) that each require different levels of precision in location specification.

\subsection{Tier 1: The Pen \& Paper Method (Mathematical Formulation)}

The yard location naming convention problem can be formulated as a constraint satisfaction problem where we must assign unique identifiers to physical locations while optimizing for operational efficiency and human comprehension.

Let us define the mathematical model:

\textbf{Sets and Parameters:}
\begin{align}
B &= \{b_1, b_2, \ldots, b_m\} \text{ (set of blocks)} \\
R_b &= \{r_1, r_2, \ldots, r_{n_b}\} \text{ (set of rows in block } b\text{)} \\
P_{b,r} &= \{p_1, p_2, \ldots, p_{k_{b,r}}\} \text{ (set of positions in row } r \text{ of block } b\text{)} \\
T_{b,r,p} &= \{t_1, t_2, \ldots, t_{h_{b,r,p}}\} \text{ (set of tiers at position } (b,r,p)\text{)}
\end{align}

\textbf{Decision Variables:}
\begin{align}
x_{b,r,p,t}^{id} &\in \{0,1\} \text{ (1 if location identifier } id \text{ is assigned to position } (b,r,p,t)\text{)} \\
s_{id} &\in \mathbb{N} \text{ (string length of identifier } id\text{)} \\
d_{id_1,id_2} &\in \mathbb{R}^+ \text{ (dissimilarity measure between identifiers } id_1, id_2\text{)}
\end{align}

\textbf{Objective Function:}
Minimize the total cognitive load while maximizing operational efficiency:
\begin{align}
\min \sum_{id} \alpha \cdot s_{id}^2 + \beta \sum_{id_1,id_2} \frac{1}{d_{id_1,id_2} + \epsilon} + \gamma \sum_{b,r,p,t} C_{access}(b,r,p,t)
\end{align}

where $\alpha$ penalizes long identifiers, $\beta$ penalizes similar identifiers, and $\gamma$ incorporates access frequency costs.

\textbf{Constraints:}
\begin{align}
\sum_{id} x_{b,r,p,t}^{id} &= 1 \quad \forall b,r,p,t \text{ (unique assignment)} \\
\sum_{b,r,p,t} x_{b,r,p,t}^{id} &\leq 1 \quad \forall id \text{ (unique identifier)} \\
s_{id} &\leq L_{max} \quad \forall id \text{ (maximum length constraint)} \\
d_{id_1,id_2} &\geq \delta_{min} \quad \forall id_1 \neq id_2 \text{ (minimum dissimilarity)}
\end{align}

\begin{figure}[htbp]
\centering
\includegraphics[width=\textwidth]{../figures-png/line6.fig1.png}
\caption{Hierarchical Yard Location Naming Structure}
\end{figure}

\textbf{Concrete Example:}

Consider a simplified yard with 2 blocks, each containing 3 rows of 4 positions with 2 tiers. We apply the constraint satisfaction formulation:

\textbf{Given Data:}
- Blocks: A, B
- Rows per block: 3
- Positions per row: 4  
- Tiers per position: 2
- Total locations: $2 \times 3 \times 4 \times 2 = 48$

\textbf{Naming Convention Selection:}
Using the format \texttt{B-RR-PP-T}, we generate identifiers systematically:

For Block A, Row 1: \texttt{A-01-01-1, A-01-01-2, A-01-02-1, A-01-02-2, A-01-03-1, A-01-03-2, A-01-04-1, A-01-04-2}

\textbf{Verification of Constraints:}
- Uniqueness: All 48 identifiers are distinct
- Length constraint: Each identifier has exactly 8 characters (within $L_{max} = 10$)
- Dissimilarity: Minimum Hamming distance between any two identifiers is 1 (satisfies $\delta_{min} = 1$)

\textbf{Cognitive Load Calculation:}
Total string length cost: $48 \times 8^2 \times \alpha = 3,072\alpha$

This mathematical foundation ensures systematic, scalable naming while optimizing for human comprehension and operational efficiency.

\subsection{Tier 2: The Classic Heuristic (Python Implementation)}

The systematic generation of optimal naming conventions can be addressed through a greedy constructive heuristic that builds naming schemes incrementally while maintaining key operational properties.

\textbf{Algorithm Theory:}

The Hierarchical Greedy Naming Algorithm constructs location identifiers by making locally optimal choices at each level of the spatial hierarchy. The algorithm prioritizes readability, uniqueness, and operational efficiency by selecting naming components that minimize confusion while maximizing information density.

\textbf{Time Complexity Analysis:}
- Preprocessing: $O(|B| \times |R| \times |P| \times |T|)$
- Identifier generation: $O(n \log n)$ where $n$ is total locations
- Conflict resolution: $O(n^2)$ in worst case
- Overall complexity: $O(n^2)$

\begin{figure}[htbp]
\centering
\includegraphics[width=\textwidth]{../figures-png/line6.fig2.png}
\caption{Hierarchical Greedy Naming Algorithm Flow}
\end{figure}

\textbf{Detailed Pseudocode:}

\lstinputlisting[language=Python]{.././codes-python/line6.py1.py}

\textbf{Concrete Example:}

Let's apply the hierarchical greedy heuristic to a container terminal with 3 blocks, varying row counts, and different position densities:

\lstinputlisting[language=Python]{.././codes-python/line6.py2.py}

\textbf{Expected Output:}
\begin{verbatim}
Naming Convention Results:
Total locations: 1896
Average identifier length: 8.00
Minimum edit distance: 1
Readability score: 95.2%

Sample Location Identifiers:
('North', 1, 1, 1) -> N-01-01-1
('North', 1, 1, 4) -> N-01-01-4
('Central', 3, 15, 2) -> C-03-15-2
('South', 6, 18, 3) -> S-06-18-3
\end{verbatim}

The greedy heuristic successfully generates 1,896 unique location identifiers with consistent 8-character length, ensuring minimum edit distance of 1 between any two identifiers and achieving 95.2\% readability score through consistent \texttt{B-RR-PP-T} formatting.

\subsection{Tier 3: The Advanced Algorithm (Metaheuristic Implementation)}

The optimal yard location naming convention problem, when considering global optimization across multiple objectives (readability, memorability, error resistance, and operational efficiency), becomes a complex combinatorial optimization challenge best addressed through evolutionary metaheuristics.

\textbf{Genetic Algorithm Theory:}

The Genetic Algorithm for Naming Convention Optimization (GANCO) evolves populations of complete naming schemes through selection, crossover, and mutation operations. Each individual in the population represents a complete assignment of identifiers to all yard locations, with fitness evaluated based on multiple criteria including cognitive load, error probability, and operational efficiency.

\textbf{Representation Scheme:}
Each chromosome encodes a complete naming convention as a vector of identifier components. For a yard with $n$ locations, the chromosome has length $4n$ (block, row, position, tier for each location).

\textbf{Fitness Function:}
\begin{align}
F(x) = w_1 \cdot f_{unique}(x) + w_2 \cdot f_{length}(x) + w_3 \cdot f_{distance}(x) + w_4 \cdot f_{pattern}(x)
\end{align}

where:
- $f_{unique}(x)$: Uniqueness penalty (0 if all unique, high penalty otherwise)
- $f_{length}(x)$: Average identifier length (shorter is better)
- $f_{distance}(x)$: Minimum edit distance between identifiers (higher is better)
- $f_{pattern}(x)$: Pattern consistency score (higher is better)

\begin{figure}[htbp]
\centering
\includegraphics[width=\textwidth]{../figures-png/line6.fig3.png}
\caption{Genetic Algorithm for Naming Convention Optimization}
\end{figure}

\textbf{Detailed Pseudocode:}

\lstinputlisting[language=Python]{.././codes-python/line6.py3.py}

\textbf{Concrete Example:}

Let's apply the genetic algorithm to optimize naming conventions for a complex multi-block terminal:

\lstinputlisting[language=Python]{.././codes-python/line6.py4.py}

\textbf{Expected Optimization Results:}
\begin{verbatim}
Total locations to optimize: 3276

Starting optimization with 3276 locations...
Generation 0: Best fitness = 2847.32
Generation 10: Best fitness = 3421.67
Generation 20: Best fitness = 3756.89
Generation 30: Best fitness = 3892.45
...
Generation 85: Best fitness = 4167.23
Generation 90: Best fitness = 4167.23

Optimization Complete!
Final fitness score: 4167.23
Generations executed: 90

Sample optimized identifiers:
Location (0, 1, 1, 1) -> N-01-01-1
Location (1, 5, 15, 3) -> C-05-15-3
Location (2, 8, 20, 1) -> S-08-20-1
Location (3, 4, 15, 2) -> W-04-15-2

Quality Metrics:
Unique identifiers: 3276/3276
Average identifier length: 8.00
Success rate: 100.0%
\end{verbatim}

The genetic algorithm successfully evolves an optimal naming convention achieving 100\% uniqueness across 3,276 locations while maintaining consistent 8-character identifiers and maximizing operational efficiency through the \texttt{B-RR-PP-T} pattern.

## Practice Questions

\textbf{Tier 1 Questions (Mathematical Formulation):}

\textbf{1.} A container terminal has 4 blocks (A, B, C, D), with Block A having 6 rows of 20 positions each, Block B having 8 rows of 24 positions each, Block C having 5 rows of 18 positions each, and Block D having 10 rows of 30 positions each. Each position can stack containers up to 4 tiers high. Formulate the constraint satisfaction problem to assign unique location identifiers with maximum length constraint of 10 characters and minimum edit distance of 2 between any two identifiers. Calculate the total number of decision variables and constraints needed.

\textbf{Expected Solution Approach:} Define binary variables $x_{b,r,p,t}^{id}$ for each location-identifier assignment. Total locations = $6 \times 20 \times 4 + 8 \times 24 \times 4 + 5 \times 18 \times 4 + 10 \times 30 \times 4 = 2,832$. Decision variables = $2,832 \times |\text{possible identifiers}|$. Constraints include uniqueness (2,832 equations) and edit distance constraints.

\textbf{2.} Given a yard with geometric constraints where containers in positions $(x,y)$ can only be accessed if positions $(x-1,y)$ and $(x,y-1)$ are accessible, formulate the accessibility constraint as part of the naming convention optimization problem. Include precedence relationships in the mathematical model.

\textbf{Expected Solution Approach:} Add accessibility variables $a_{x,y} \in \{0,1\}$ and precedence constraints: $a_{x,y} \leq a_{x-1,y}$ and $a_{x,y} \leq a_{x,y-1}$ for all valid $(x,y)$. Link naming decisions to accessibility through additional constraints.

\textbf{Tier 2 Questions (Heuristic Implementation):}

\textbf{3.} Implement a greedy heuristic that assigns location identifiers by prioritizing high-traffic areas first. Given the following access frequency data: North Block (1,200 moves/day), Central Block (2,800 moves/day), South Block (800 moves/day), and West Block (450 moves/day), design an algorithm that assigns shorter identifiers to frequently accessed locations while maintaining systematic ordering.

\textbf{Expected Solution Approach:} Sort blocks by frequency, assign single-character block IDs to highest traffic areas (C=Central, N=North, S=South, W=West). Implement priority queue for position assignment within each block. Target average identifier length $\leq 7.5$ characters.

\textbf{4.} Design a constructive heuristic that handles naming convention evolution when new blocks are added to an existing terminal. The algorithm must maintain backward compatibility with existing identifiers while integrating new areas seamlessly. Test with a scenario where blocks A-F exist, and blocks G-K are being added with different geometric layouts.

\textbf{Expected Solution Approach:} Implement identifier versioning scheme, reserve namespace for expansion, use hierarchical prefix allocation. Ensure no conflicts with existing 2,400 identifiers while adding 1,800 new locations.

\textbf{Tier 3 Questions (Metaheuristic Optimization):}

\textbf{5.} Configure a genetic algorithm to optimize naming conventions for a terminal experiencing seasonal capacity changes. During peak season, temporary overflow areas T1-T4 are activated (adding 1,200 positions), while off-peak operations use only permanent areas. The algorithm must optimize for both scenarios simultaneously with weighted objectives: 60\% permanent operations, 40\% peak operations.

\textbf{Expected Solution Approach:} Design dual-fitness evaluation function, implement scenario-based crossover operators. Population size = 80, generations = 150, crossover rate = 0.75, mutation rate = 0.12. Target combined fitness $> 4,500$.

\textbf{6.} Develop a multi-objective genetic algorithm that optimizes naming conventions considering three conflicting objectives: (1) minimize cognitive load for human operators, (2) maximize machine readability for automated systems, and (3) minimize identifier length for label printing costs. Use Pareto optimization with the following constraints: 5,500 total locations, identifier length $\leq 12$ characters, minimum edit distance $\geq 2$.

\textbf{Expected Solution Approach:} Implement NSGA-II or MOEA/D algorithm. Define cognitive load as pattern complexity measure, machine readability as parsing efficiency score. Generate Pareto front with 20-30 non-dominated solutions. Report hypervolume indicator and convergence metrics.

\end{document}
